% The outline of the whole paper (PI 2026-11-24, LEDGER RW-10): one sentence per paragraph, under the section plan
% and word budgets of docs/STYLE.md section 1, with the place of each of the 5 figures and 4 tables.
\newcommand{\osec}[2]{\par\smallskip\noindent\textbf{#1}\quad\textit{#2}\par}
\newcommand{\op}[1]{\par\noindent\hangindent=1.4em\hangafter=1\makebox[1.4em][l]{#1.}}
{\fontsize{8.6}{10}\selectfont\setlength{\parskip}{0pt}
\noindent{\normalsize\textbf{Outline: one sentence per paragraph}}\par
\osec{1 Introduction}{about 600 words}
\op{1} Particle-based networks have made jet tagging central to LHC analyses.
\op{2} Jet foundation models are pretrained by classification into many fine classes, as Sophon is on JetClass-II.
\op{3} Whether fine labels improve the model is untested: published comparisons change the jets with the labels.
\op{4} We relabel the same jets with four nested label sets, random groupings and a mass regression, and change nothing else.
\op{5} Every model is evaluated frozen and fine-tuned, on JetClass-II, JetClass, top tagging and quark--gluon discrimination.
\op{6} The findings, in the order of the abstract.
\osec{2 Related work}{about 350 words, no numbers from other papers}
\op{1} Supervised classification is the reference objective for jets; no study has varied its label set with the jets fixed.
\op{2} In image classification finer labels help more for frozen features than after fine-tuning, depending on the grouping.
\op{3} Auxiliary targets such as the jet mass, and transfer from frozen features, have not been studied against the label set.
\osec{3 Methods}{about 1,400 words; Table 1 and Fig.~1 (design) in 3.2}
\op{1} \textit{3.1 Data and model.} JetClass-II, its \Kfull{} classes and the selection of the jets.
\op{2} The Sophon network and the frozen features that probes and transfer learning read.
\op{3} Every model sees the same jets in the same order from the same initial weights.
\op{4} \textit{3.2 Label sets.} The four nested label sets and the distinctions each merges (Table~1, Fig.~1).
\op{5} A merged class scores the sum of the classes it merges (equation), which fixes the discriminants a label set can form.
\op{6} \textit{3.3 Controls.} Random groupings keep the class count and class sizes of the \Kcoarse-class set but not its physics.
\op{7} The mass regression adds one output and one loss term (equation) to the \Kres- and \Kcoarse-class models.
\op{8} \textit{3.4 Evaluation.} Linear probes on the frozen features: JetClass-II classes, class pairs, benchmarks; one sentence on missing inputs.
\op{9} Fine-tuning with one schedule for every model, and the same network trained from scratch as reference.
\op{10} Accuracy, AUC and Rej$_{30\%}$ as in ParT, and the mass resolution (equation).
\op{11} Sophon's $X\to b\bar b$ and $X\to bs$ benchmark by transfer learning on the frozen features.
\op{12} \textit{3.5 Uncertainties.} Mean and standard deviation over pretraining runs; ``significant'' is beyond three standard deviations (run spreads and test-sample bootstrap).
\osec{4 Results}{about 2,600 words}
\op{1} \textit{4.1 The pretraining classes (Table~2, Fig.~2: frozen-feature summary).} Linear-probe accuracy on the \Kfull{} classes falls with coarser labels, from \VtwoAccFullPct\% to \VtwoAccCoarsePct\%.
\op{2} The \Kres-class model loses only on the QCD classes it merges.
\op{3} Random groupings lose far less than the \Kcoarse-class model, so which classes are merged matters.
\op{4} The untrained network and Sophon (released) set the scale; one sentence on the MLP probe.
\op{5} \textit{4.2 Class pairs (Fig.~3: two pairs, label sets and random groupings).} A pair that differs only in a merged flavor ($X\to b\bar b$ against $X\to c\bar c$) becomes much harder to separate.
\op{6} A pair loses little until the label set merges its flavor.
\op{7} A discriminant summed from fine classes therefore needs those classes among the pretraining labels.
\op{8} \textit{4.3 Fine-tuning (Table~3, Fig.~4: fine-tuning against training size).} JetClass: \pending{pending the fine-tuned JetClass cells}.
\op{9} Top tagging: no significant difference between label sets; one sentence on the run-to-run spread.
\op{10} Quark--gluon discrimination: no significant difference; the same sentence.
\op{11} Training size: \pending{pending the cells below the full sets}.
\op{12} Pretrained models against the network trained from scratch.
\op{13} \textit{4.4 Jet mass (Table~4).} The mass regression sharpens the mass resolution obtained from the frozen features.
\op{14} It has no significant cost in classification.
\op{15} \textit{4.5 Resonance benchmark (Fig.~5: $X\to b\bar b$ and $X\to bs$).} Transfer learning on the frozen features loses signal efficiency for both resonances at \Kcoarse{} classes.
\op{16} On $X\to bs$ the \Kfull-class model stays well above the same classifier trained from scratch.
\osec{5 Discussion and conclusion}{about 600 words, prose}
\op{1} A label set should keep the flavor distinctions that analyses read from frozen features or summed scores.
\op{2} After fine-tuning: \pending{pending the JetClass result}.
\op{3} Limitations in three sentences: one architecture and dataset; simulation only; one step merges several distinctions at once.
\op{4} Conclusion.
\osec{Appendix}{at most 15 pages}
\op{A} Per-class recall, other class pairs, probes against training size, mass against $p_{\mathrm T}$, validation curves, published taggers, Herwig, AUC columns, flavor-boundary groupings, MLP probes, class map, inputs, learning rates.
\par}
